\section{Entire original Laurent coefficients of original pulled-module sections}

\begin{theorem}\label{thm:fc-scheme-module-laurent-coefficients}
\lean{EnvelopeProperty.SchemeModuleLaurentCoefficients.pulledSections,EnvelopeProperty.SchemeModuleLaurentCoefficients.coefficients}
\leanok
\uses{thm:fc-flat-module-global-sections,thm:laurent-module-actual-coefficients}
Let $R$ be an original commutative ring, $p:X\to\operatorname{Spec}R$
an original compact separated scheme, and $M$ any original
quasicoherent module. Put $S=R[T,T^{-1}]$, $X_S=X\times_R\operatorname{Spec}S$
and retain the original projections $g$ and $q$. There is a canonical
bijection of the entire original section modules
\[
 \Gamma(X_S,g^*M)\simeq\mathbb Z\to_0\Gamma(X,M).
\]
The original coefficient module retains its literal original $R$ action;
the pulled module retains its literal original $S$ action through $q$.
No finite-dimensionality, freeness of $M$, chosen Laurent expansion,
or global-section comparison is assumed.
\end{theorem}
\begin{proof}
The original Laurent algebra has its actual integer-indexed monomial
basis over $R$, and is therefore flat over $R$. Invert the preceding
whole original flat module H0 equivalence with this actual algebra.
Compose it with the already proved entire Laurent tensor-module
coefficient equivalence. This gives the stated bijection on the
literal original modules and determines every coefficient uniquely.
The scalar actions are the same actual actions in those two original
producers. Only the underlying bijection is asserted here; no further
identification of scalar structures is needed.
\end{proof}

\begin{theorem}\label{thm:fc-scheme-module-laurent-coefficient-formulas}
\lean{EnvelopeProperty.SchemeModuleLaurentCoefficients.coefficients_tensorEquiv,EnvelopeProperty.SchemeModuleLaurentCoefficients.coefficients_pure,EnvelopeProperty.SchemeModuleLaurentCoefficients.coefficients_monomial,EnvelopeProperty.SchemeModuleLaurentCoefficients.coefficients_symm_tensorEquiv,EnvelopeProperty.SchemeModuleLaurentCoefficients.coefficients_symm,EnvelopeProperty.SchemeModuleLaurentCoefficients.eq_sum_monomial}
\leanok
\uses{thm:fc-scheme-module-laurent-coefficients,thm:laurent-module-actual-coefficients}
For every original Laurent polynomial $F$, every original section $m$
and every integer $i$, the coefficients of
$q^*(F)\eta_g(m)$ at $i$ are $F_i m$. In particular the whole coefficient
family of $q^*(T^i)\eta_g(m)$ is the original singleton at $i$ with
value $m$. Conversely every original finite coefficient family $c$
reconstructs the original section
\[
 \sum_{i\in\operatorname{supp}(c)}q^*(T^i)\eta_g(c_i).
\]
Thus every original pulled-module section has this actual finite
Laurent expansion, with its uniquely determined original coefficients.
\end{theorem}
\begin{proof}
Apply the actual H0 pure-tensor formula to $F\otimes m$, then use the
original Laurent tensor coefficient formula. Invertibility gives the
monomial singleton and inverse assertions. The already proved tensor
inverse is its actual finite sum of monomial pure tensors. The whole
H0 map is linear, so it carries that finite sum to exactly the stated
sum of original coefficient times original pullback-unit sections.
Applying the two inverse identities gives the expansion of every
original section. No expansion or coefficient compatibility is
supplied as a premise.
\end{proof}

\begin{definition}\label{def:fc-scheme-module-laurent-finite-coefficient-family}
\lean{EnvelopeProperty.SchemeModuleLaurentCoefficients.terms,EnvelopeProperty.SchemeModuleLaurentCoefficients.terms_eq_tensor_terms}
\leanok
\uses{thm:fc-scheme-module-laurent-coefficients,thm:laurent-module-actual-finite-family}
For any original finite family $G$ of sections of $g^*M$ on $X_S$,
form the union, over its original sections, of their actual coefficient
supports. The pairs $(c_i,i)$ form one finite family of original
sections of $M$ and original integer weights. This entire family is
chosen from $G$ before any point, residue map, valuation or radius.
It is exactly the finite pair family obtained by sending $G$ through
the inverse original H0 map and applying the original Laurent tensor
coefficient-family construction.
\end{definition}
\begin{proof}
Each actual coefficient family is finitely supported. A finite union
of their finite support images is finite. The two definitions agree
because the geometric coefficient map is the original Laurent tensor
coefficient map after the inverse whole H0 equivalence. Taking the
image of $G$ under that inverse map and distributing the finite union
proves the exact pair-family equality. The construction has no point,
residue, valuation or radius input.
\end{proof}
